% REPORT-SHARD: 03_ising_audit_and_numerics
\section{The Ising audit and the finite fermion evidence}

\subsection{A stronger local limit identification}

The initial proof reconstructed the Ising interval net from sampled smooth
even field products. A second skeptical audit established an upper inclusion
for every uniformly bounded local lattice sequence with a represented weak
limit, under the same sharp CAR refinements. This removes a possible concern
that the construction selected only a convenient subset of local limits.

\textbf{Lemma A2 [Lamport proof].}
\textbf{Assume} $B_M$ are even lattice observables supported in an interval
$I$, with $\sup_M\|B_M\|<\infty$, and their represented images converge
weakly to $B$ on the NS even-vacuum space. \textbf{Prove} $B\in\mathcal A(I)$.
\begin{description}
\item[$\langle1\rangle1$.] Let $D$ be an even smooth field polynomial
supported compactly in the complementary interval $I'$. Its sampled
approximants $D_M$ converge in norm by LL-1/LL-2.
\item[$\langle1\rangle2$.] Finite disjoint spatial support gives
$[B_M,D_M]=0$. Uniform boundedness and norm convergence of $D_M$ imply
$[B_M,D]\to0$ in norm.
\item[$\langle1\rangle3$.] Weak convergence of $B_M$ then gives
$[B,D]=0$. The complementary sampled polynomials generate
$\mathcal A(I')$, so $B\in\mathcal A(I')'$.
\item[$\langle1\rangle4$.] Haag duality of the identified Ising vacuum
net gives $\mathcal A(I')'=\mathcal A(I)$. QED by
$\langle1\rangle1$--$\langle1\rangle3$ and the sourced duality theorem.
\end{description}
The sampled generators supply the opposite inclusion. Thus the von Neumann
algebra generated by all such bounded local limits is the same interval
algebra. This statement does not make arbitrary weak convergence
multiplicative, or replace the strict OAR algebra maps by a sequence quotient.

\sources{\path{cft_machine/reviews/second_pass_ising.md}, A1/A2;
\path{cft_machine/reviews/oar_second_pass_referee.md}; KL source lines
403--407 (Haag duality); OS source lines 1451--1558 (operator convergence).}

\subsection{What happens to the ultraviolet doubling}

The local sine Hamiltonian has additional finite low-energy excitations near
the edge of the Brillouin zone. For
$r_M=M/2-1/2$, its positive one-particle energy is
\[
 \frac{\sin(\epsilon_Mr_M)}{\epsilon_M}
   =\frac{M}{2\pi}\sin(\pi/M)\longrightarrow\frac12.
\]
An even pair at the top two momenta similarly has limiting energy $2$.
These are not removed simply by passing to even observables.

However, the chosen sharp refinement preserves a fixed momentum label.
The labels of these edge excitations keep changing with $M$; their embedded
vectors have no norm limit. The proved result is therefore the specified
chiral OAR limit, not convergence of the entire sequence of low-energy
spectra. The source explicitly permits sharp, spatially nonlocal scaling
maps. Its separate wavelet-locality proposition was not misattributed to
this sharp branch.

\subsection{Why the modified KS prescription matters}

For the complex chiral sea, the Schwinger cocycle is
\[
 c_S(A,B)=\Tr\{SA(1-S)BS-SB(1-S)AS\}.
\]
With the continuum shifts $B_ne_r=(r-n/2)e_{r-n}$ it reduces, for $n>0$,
to a finite sum over the $n$ crossings of the Fermi surface:
\[
 \sum_{j=0}^{n-1}\left(j+\frac{1-n}{2}\right)^2
   =\frac{n(n^2-1)}{12}.
\]
The self-dual Majorana value is half of this. These are normal-ordering
terms in second quantization, not central terms of bare finite matrices.

\begin{table}[h]
\centering
\begin{tabular}{@{}rcc@{}}
\toprule
$M$ & Complex cocycle, $n=2$ & Complex cocycle, $n=3$\\\midrule
8 & 0.3459321593 & 0.7748288164\\
16 & 0.4567488187 & 1.5926816948\\
32 & 0.4888651513 & 1.8903265978\\
Exact continuum & $1/2$ & $2$\\\bottomrule
\end{tabular}
\caption{Source-modified finite KS modes. Exact target values are derived
from the crossing sum, not fitted to the samples.}
\end{table}
Periodic momentum wrap adds an ultraviolet crossing and cancels the cocycle
in this diagnostic. The source's characteristic-function correction removes
that spurious wrap; a deliberate mutation of the implementation is rejected.

On a fixed mode core the scalar coefficient error has the analytic bound
\[
 \left|a_n^M(r)-(r-n/2)\right|\le\epsilon_M^2
 \left(\frac{|r-n/2|^3}{6}+\frac{n^2|r-n/2|}{16}\right).
\]
Adding the bounds at adjacent scales gives a summable refinement defect.
The run explicitly marks $M=8,n=3$ as outside its chosen shifted-core
hypothesis; a displayed bound is not silently applied beyond its domain.

\subsection{Actual finite Dirac states, rather than a substituted sea}

The two-component Dirac benchmark has finite negative-energy projections
$P_Q(r)=(1-h_Q(r)/E_Q(r))/2$. Their row pullbacks change with $Q$.
For fixed coarse momentum set $\Gamma_M$,
\[
 \|P_{M,Q}-P_{M,\infty}\|
   =\sin\!\left(\frac{\epsilon_Q\max_{r\in\Gamma_M}|r|}{4}\right).
\]
The proof then converts covariance convergence into a genuine state-norm
bound. Each retained two-mode block is a pure occupied-orbital state;
its density-matrix distance is twice the projection norm. The product
state is handled by telescoping tensor factors, giving
\[
 \|\omega_{M,Q}-\omega_{M,\infty}\|
   \le\frac{\epsilon_Q}{2}\sum_{r\in\Gamma_M}|r|.
\]
Adjacent fine-scale drift is bounded by half this coefficient. The dyadic
tail is computable and summable for each fixed row. It need not be uniform
in the coarse row, and may be a loose bound at small cutoffs.

\sources{\path{cft_machine/benchmarks/fermion/PROOF.md};
\path{cft_machine/benchmarks/fermion/WILSON_PROOF.md};
\path{cft_machine/benchmarks/fermion/runs/2026-09-22/results.toml}.
The source convergence theorem supplies the analytic continuum result;
the runs verify the implemented formulas and independent finite oracles.}
